\documentclass[border=5pt,svgnames,tikz]{standalone}
\usepackage{fontspec}
\usepackage{ctex}
\setmainfont{Source Serif 4}
\setsansfont{Source Sans 3}
\setmonofont{Source Code Pro}
\setCJKmainfont[BoldFont=Source Han Sans CN]{Source Han Sans CN}
\usetikzlibrary{positioning, calc}

\begin{document}
\begin{tikzpicture}[
    node distance=0pt,
    box/.style={draw, rectangle, minimum size=0.6cm, inner sep=0pt, anchor=center, font=\sffamily},
    label text/.style={font=\sffamily\small, align=center}
]

\definecolor{custompurple}{HTML}{B266FF}

% Define the base sequence
% 1:U, 2:U, 3:F, 4:U, 5:F, 6:U, 7:U, 8:U, 9:U, 10:F, 11:F, 12:U, 13:U, 14:U, 15:F, 16:F
% U = 1/0, F = 0/1
% 1, 1, 0, 1, 0, 1, 1, 1, 1, 0, 0, 1, 1, 1, 0, 0
\def\baseseq{1,1,0,1,0,1,1,1,1,0,0,1,1,1,0,0}

% Macro to draw a sub-diagram
% Arguments: Name, Position, OverrideIndex, OverrideType, Label, Description
% OverrideType:
% 0: None (Consistent)
% 1: Missing (0/0 at index)
% 2: DupFree (0/2 at index)
% 3: DupData (2/0 at index)

\newcommand{\drawdiagram}[6]{
    \begin{scope}[shift={#2}, local bounding box=#1]
        \node[anchor=east, font=\sffamily\small] at (-0.5, 0.3) {使用表};
        \node[anchor=east, font=\sffamily\small] at (-0.5, -0.3) {空闲表};

        \foreach \val [count=\i] in \baseseq {
            \pgfmathsetmacro{\x}{(\i-1)*0.6}

            % Default values
            \def\topval{\val}
            \pgfmathsetmacro{\botval}{1-\val}
            \def\topcolor{black}
            \def\botcolor{black}

            % Overrides
            \ifnum #3=\i
                \ifnum #4=1 % Missing
                    \def\topval{0}
                    \def\botval{0}
                    \def\botcolor{red}
                \fi
                \ifnum #4=2 % DupFree
                    \def\topval{0}
                    \def\botval{2}
                    \def\botcolor{red}
                \fi
                \ifnum #4=3 % DupData
                    \def\topval{2}
                    \def\botval{0}
                    \def\topcolor{red}
                \fi
            \fi

            \node[box] (T\i) at (\x, 0.3) {\textcolor{\topcolor}{\pgfmathprintnumber[int trunc]{\topval}}};
            \node[box] (B\i) at (\x, -0.3) {\textcolor{\botcolor}{\pgfmathprintnumber[int trunc]{\botval}}};
        }

        % Label a), b), c), d) and Description
        % Center them below the diagram
        % The diagram width is 16 * 0.6 = 9.6
        % Center is at 9.6 / 2 - 0.3 (half box) = 4.5
        % Actually x goes from 0 to 15*0.6 = 9.0. Center is 4.5.

        \node[below=0.3 of B8, anchor=center, xshift=0.3cm] (lbl) {\small #5 \quad \textcolor{custompurple}{#6}};
    \end{scope}
}

% a) Consistent
\drawdiagram{A}{(0,0)}{0}{0}{a)}{一致}

% b) Missing (Index 3)
\drawdiagram{B}{(11.5,0)}{3}{1}{b)}{块丢失}

% c) DupFree (Index 5)
\drawdiagram{C}{(0,-2)}{5}{2}{c)}{重复空闲块}

% d) DupData (Index 6)
\drawdiagram{D}{(11.5,-2)}{6}{3}{d)}{重复数据块}

\end{tikzpicture}
\end{document}
